%%% FIELD finite-quantile-statement
For each fixed \(N\) and allowed observed design, \(\Lambda_N^{\mathrm{phys}}\) is compact.  The law of
\[
d_{H_N}(G_N+\lambda q_N(p),\mathcal C_N)^2
\]
is atomless and has a strictly increasing distribution function on \((0,\infty)\); its \((1-\alpha)\)-quantile is continuous in \((\lambda,p)\).  Hence the maximum defining \(c_{N,\alpha}^{\mathrm{att}}\) is attained and
\[
 \inf_{(\lambda,p)\in\Lambda_N^{\mathrm{phys}}}
 \Pr\{d_{H_N}(G_N+\lambda q_N(p),\mathcal C_N)^2
       \le c_{N,\alpha}^{\mathrm{att}}\}=1-\alpha.                           \tag{20}
\]
Moreover,
\[
 0<c_{N,\alpha}^{\mathrm{att}}\le\chi^2_{2,1-\alpha},                       \tag{21}
\]
and both bounds are uniform over the allowed designs after compactification.
%%% FIELD finite-quantile-proof
Fix \(N\) and an allowed design.  Let \(\Theta_N\) be the subset of
\[
[-A_{\max},A_{\max}]\times[p_{\min},p_{\max}]
   \times\mathcal K_{\mathrm{nuis}}
\]
on which the determinant-margin inequality holds at every \(\omega\in\mathbb R^d\).  The displayed product is compact by the amplitude, exponent, and nuisance-compactness assumptions.  For each fixed \(\omega\), the determinant-margin inequality is a closed condition in the coordinates of this product, because the three spectra are continuous functions of those coordinates.  Hence \(\Theta_N\), an intersection of closed subsets of a compact set, is compact.  The ideal stencil contains finitely many fixed distances, so its power-law moment is a finite linear combination of functions \(r^{2p}\).  Thus \(P_\star(p)\), and hence \(K_d(p)\), is continuous on the exponent interval; positivity is part of their definitions.  Since \(h_N,c_X,c_W>0\), the map
\[
(A,p,c_X,c_W,\alpha_X,\alpha_W,\alpha_C)
 \longmapsto
 \left(\frac{A K_d(p)\sqrt N h_N^{2p-1}}{\sqrt{c_Xc_W}},p\right)             \tag{22}
\]
is continuous.  By the definition of the physically attainable set, its image is \(\Lambda_N^{\mathrm{phys}}\).  An allowed model supplies at least one legal tuple, so this image is nonempty.  It is therefore compact.

Set
\[
S_{N,\lambda,p}=d_{H_N}(G_N+\lambda q_N(p),\mathcal C_N)^2.
\]
Write \(F_{N,\lambda,p}\) for the distribution function of \(S_{N,\lambda,p}\).  The profile-geometry compactification lemma gives
\[
F_{N,\lambda,p}(0)=0,
\qquad F_{N,\lambda,p}\text{ is continuous on }[0,\infty),
\qquad F_{N,\lambda,p}\text{ is strictly increasing on }(0,\infty).         \tag{23}
\]
In particular the quantile
\[
k_N(\lambda,p)=\inf\{x\ge0:F_{N,\lambda,p}(x)\ge1-\alpha\}                 \tag{24}
\]
is positive, unique, and satisfies
\[
F_{N,\lambda,p}\{k_N(\lambda,p)\}=1-\alpha.                                \tag{25}
\]
If \((\lambda_r,p_r)\to(\lambda,p)\), use the same Gaussian vector \(G_N\) for all indices.  Continuity of \(q_N\) and the fact that distance to a fixed closed set is continuous give
\[
S_{N,\lambda_r,p_r}\longrightarrow S_{N,\lambda,p}
\quad\text{almost surely}.                                                    \tag{26}
\]
Equation (23) and (26) imply pointwise convergence of the corresponding distribution functions.  If the quantiles in (24) failed to converge to \(k_N(\lambda,p)\), a subsequence would lie on one side of a fixed intermediate point.  Evaluating the distribution functions at that point and using (23) would contradict (24)--(25).  Thus \(k_N\) is continuous.  Compactness of \(\Lambda_N^{\mathrm{phys}}\) now supplies a maximizer \((\lambda_*,p_*)\), and
\[
c_{N,\alpha}^{\mathrm{att}}=k_N(\lambda_*,p_*).                              \tag{27}
\]
For every \((\lambda,p)\) in the physical set, (25) and (27) give
\[
F_{N,\lambda,p}(c_{N,\alpha}^{\mathrm{att}})\ge1-\alpha,
\]
while equality holds at \((\lambda_*,p_*)\).  This proves (20).

For the upper bound, the cone contains the line \(\{a q_N(p):a\in\mathbb R\}\).  Therefore
\[
\begin{aligned}
S_{N,\lambda,p}
&\le \inf_{a\in\mathbb R}
 \|G_N+\lambda q_N(p)-a q_N(p)\|_{H_N}^2\\
&=\inf_{b\in\mathbb R}\|G_N-bq_N(p)\|_{H_N}^2,                              \tag{28}
\end{aligned}
\]
where \(\|x\|_{H_N}^2=x^{\mathsf T}H_N^{-1}x\).  After whitening, the last quantity is the squared norm of the projection of a standard three-dimensional Gaussian vector onto the two-dimensional orthocomplement of \(H_N^{-1/2}q_N(p)\).  It has the \(\chi^2_2\) law.  Equations (24) and (28) prove the upper inequality in (21).

Finally, the profile-geometry lemma extends the quantile map continuously to a sequentially compact family containing all bounded-amplitude, tangent-plane, and endpoint-half-plane limits, including limits of nonconvergent designs.  At every point of that compactification the limiting squared distance has probability zero at zero by the same lemma.  Since \(1-\alpha>0\), its \((1-\alpha)\)-quantile is strictly positive.  A continuous positive function on a compact set has a positive minimum; call it \(c_0\).  This proves the lower inequality in (21) and the asserted uniform continuity.
%%% FIELD exact-calibration-statement
For each \(N\) and allowed observed design, let \(c_{N,\alpha}^{\mathrm{att}}\) be the exact attained critical value over \(\Lambda_N^{\mathrm{phys}}\), and let
\[
C_N^{\mathrm{att}}=\{b\in\mathcal B:T_N(b)\le c_{N,\alpha}^{\mathrm{att}}\}.
\]
The Gaussian reference quantile is uniformly continuous on the bounded-amplitude, tangent-plane, and endpoint-half-plane compactification, and
\[
\inf_{(\lambda,p)\in\Lambda_N^{\mathrm{phys}}}
\Pr\!\left\{d_{H_N}(G_N+\lambda q_N(p),\mathcal C_N)^2
\le c_{N,\alpha}^{\mathrm{att}}\right\}=1-\alpha.                           \tag{29}
\]
For the actual observed-law experiment,
\[
\lim_{N\to\infty}\inf_{(P_N,\mathcal S_N)\in\mathcal M}
P_N\{\beta\in C_N^{\mathrm{att}}\}=1-\alpha,                               \tag{30}
\]
and there is a finite constant \(C\), depending only on the fixed model and level constants, such that
\[
\operatorname{diam}(C_N^{\mathrm{att}})\le C N^{-1/2}\quad P_N\text{-almost surely}                              \tag{31}
\]
uniformly over the model class.  The conclusions include zero amplitude, bounded and diverging local amplitudes, exponent endpoints, the one-dimensional logarithmic elbow window, and nonconvergent allowed-design sequences.  Equation (30) is an asymptotic observed-law statement; at fixed \(N\), (29) is exact only for the Gaussian reference statistic.  The critical value is mathematically defined and attained, but no terminating certified numerical approximation to it is asserted.
%%% FIELD exact-calibration-proof
Fix an allowed model and design, and abbreviate
\[
E_N=\sqrt{\frac N\gamma}\{m_N-\beta\mathbf1-B_Nq_N(p)\},
\qquad
\lambda_N=\frac{\sqrt N B_N}{\sqrt\gamma},
\qquad
\Phi_{N,\lambda,p}(e)
=d_{H_N}(e+\lambda q_N(p),\mathcal C_N)^2.                                   \tag{32}
\]
The nuisance mean \(\lambda q_N(p)\) belongs to \(\mathcal C_N\).  Because \(\mathcal C_N\) is a cone and \(m_N-\beta\mathbf1=\sqrt{\gamma/N}\{E_N+\lambda_Nq_N(p)\}\), homogeneity of distance gives the exact identity
\[
T_N(\beta)
=\frac{\gamma}{\widehat\gamma_N}
 \Phi_{N,\lambda_N,p}(E_N).                                                   \tag{33}
\]

We first transfer the reference distribution uniformly.  The uniform-experiment theorem, together with the upper and lower bounds on \(\gamma=c_W/c_X\) from nuisance compactness, yields
\[
\sup_{(P_N,\mathcal S_N)\in\mathcal M}
d_{\mathrm{BL}}\{\mathcal L_{P_N}(E_N),N_3(0,H_N)\}\longrightarrow0.        \tag{34}
\]
The eigenvalues of \(H_N\) are bounded above and away from zero by the same theorem.  On the Euclidean ball \(\{e:\|e\|\le R\}\), distance to a closed set is one-Lipschitz in its defining metric and comparison with the cone point \(\lambda q_N(p)\) gives
\[
d_{H_N}(e+\lambda q_N(p),\mathcal C_N)\le\|e\|_{H_N}.
\]
Consequently, for \(\|e\|,\|e'\|\le R\),
\[
|\Phi_{N,\lambda,p}(e)-\Phi_{N,\lambda,p}(e')|
\le \{\|e\|_{H_N}+\|e'\|_{H_N}\}\|e-e'\|_{H_N}
\le L_R\|e-e'\|,                                                            \tag{35}
\]
where \(L_R\) is independent of the design, \(\lambda\), and \(p\).  Equation (34) and a bounded Lipschitz cutoff imply uniform tightness of \(E_N\), because the centered Gaussian reference vectors are uniformly tight under the eigenvalue bound.  Smooth upper and lower approximations to the indicator of \(\{\Phi_{N,\lambda,p}\le x\}\), first on the ball in (35) and then outside it by uniform tightness, therefore show for every finite \(M\) that
\[
\begin{split}
\sup_{(P_N,\mathcal S_N)\in\mathcal M}\sup_{0\le x\le M}
\big|P_N\{\Phi_{N,\lambda_N,p}(E_N)\le x\}
{}-{}\Pr\{\Phi_{N,\lambda_N,p}(G_N)\le x\}\big|\longrightarrow0.          \tag{36}
\end{split}
\]
Indeed, for a smoothing band of width \(\varepsilon\), the bounded-Lipschitz error tends to zero by (34)--(35), and the remaining reference probability is bounded by
\[
\sup_{N,\mathcal S_N,\lambda,p}\sup_{0\le x\le M}
\Pr\{|\Phi_{N,\lambda,p}(G_N)-x|\le\varepsilon\},
\]
which tends to zero as \(\varepsilon\downarrow0\) by the uniform anti-concentration relation in the profile-geometry compactification lemma.  This proves (36) without requiring the design sequence to converge.

Let \(R_N=\gamma/\widehat\gamma_N\).  Feasible scale consistency in the uniform-experiment theorem gives \(R_N\to1\) in probability uniformly over the model.  The finite-design quantile lemma gives constants
\[
0<c_0\le c_{N,\alpha}^{\mathrm{att}}\le M_\alpha:=\chi^2_{2,1-\alpha}.        \tag{37}
\]
On \(\{|R_N-1|\le\eta\}\), with \(0<\eta<1/2\),
\[
\{\Phi_{N,\lambda_N,p}(E_N)\le c_{N,\alpha}^{\mathrm{att}}/(1+\eta)\}
\subseteq\{T_N(\beta)\le c_{N,\alpha}^{\mathrm{att}}\}
\subseteq
\{\Phi_{N,\lambda_N,p}(E_N)\le c_{N,\alpha}^{\mathrm{att}}/(1-\eta)\}.     \tag{38}
\]
The two thresholds in (38) differ from \(c_{N,\alpha}^{\mathrm{att}}\) by at most \(2M_\alpha\eta\).  Apply (36), then the same uniform anti-concentration bound, and finally let \(\eta\downarrow0\).  Uniform scale consistency removes the complement of the event in (38), giving
\[
\sup_{(P_N,\mathcal S_N)\in\mathcal M}
\left|P_N\{T_N(\beta)\le c_{N,\alpha}^{\mathrm{att}}\}
-\Pr\{\Phi_{N,\lambda_N,p}(G_N)\le c_{N,\alpha}^{\mathrm{att}}\}\right|
\longrightarrow0.                                                            \tag{39}
\]

For each fixed allowed design, the finite-design attainable-quantile lemma proves (29).  Every observed-law parameter maps, by the observation-amplitude definition, to a pair in \(\Lambda_N^{\mathrm{phys}}\).  Conversely, the definition of \(\Lambda_N^{\mathrm{phys}}\) makes each of its points the image of a legal spectral tuple, so the maximizing pair in (27) is available to the infimum over the model class.  Hence the infimum of the reference probability on the right side of (39) is exactly \(1-\alpha\) for every \(N\) and design.  Taking the uniform infimum in (39) proves both inequalities needed for (30).  Notice that the equality for the observed law follows only after the approximation error in (39) tends to zero; it is not a finite-sample observed-law equality.

It remains to prove the pathwise diameter bound.  Take \(b_1,b_2\in C_N^{\mathrm{att}}\).  For arbitrary \(\varepsilon>0\), choose \(z_i\in\mathcal C_N\) such that, with \(r_i=m_N-b_i\mathbf1-z_i\),
\[
\|r_i\|_{H_N}\le
\sqrt{\widehat\gamma_Nc_{N,\alpha}^{\mathrm{att}}/N}+\varepsilon,
\qquad i=1,2.                                                                \tag{40}
\]
The moment-separation lemma, the identity
\[
r_2-r_1=(b_1-b_2)\mathbf1+z_1-z_2,
\]
and the upper eigenvalue bound yield
\[
\begin{aligned}
c_\star|b_1-b_2|
&\le\|(b_1-b_2)\mathbf1+z_1-z_2\|\\
&\le\sqrt{\lambda_{\max}(H_N)}
   \{\|r_1\|_{H_N}+\|r_2\|_{H_N}\}.                                      \tag{41}
\end{aligned}
\]
The deterministic clipping in \(\widehat\gamma_N\) gives
\[
0<\widehat\gamma_N\le4\overline c_W/\underline c_X.                        \tag{42}
\]
Letting \(\varepsilon\downarrow0\) in (40)--(41), and using (37), (42), and the uniform eigenvalue bound, proves
\[
|b_1-b_2|\le
\frac{2}{c_\star}
\sqrt{\lambda_{\max}(H_N)\widehat\gamma_N
      c_{N,\alpha}^{\mathrm{att}}/N}
\le C N^{-1/2}.                                                               \tag{43}
\]
Taking the supremum over \(b_1,b_2\) proves (31), with the stated convention for an empty set.  The denominator control in the uniform-experiment theorem makes all local ratios well defined with probability one; defining the procedure arbitrarily on the null exceptional event preserves (31) almost surely.

Finally, the attainable-regimes lemma enumerates zero, bounded, divergent, endpoint, and one-dimensional logarithmic-elbow sequences in the physical set.  The profile-geometry lemma compactifies their tangent-plane and endpoint-half-plane limits and treats nonconvergent designs.  Since (36)--(43) are uniform over that entire compactification, no additional case-dependent assumption is used.  Neither lemma supplies an effective two-sided representation or convergence modulus for a terminating certified continuum optimization, so the proof establishes exact mathematical calibration but makes no certified-computability claim.
%%% FIELD calibration-handle-correction
For each \(N\) and observed design, form the physically attainable set \(\Lambda_N^{\mathrm{phys}}\), define
\[
c_{N,\alpha}^{\mathrm{att}}
=\max_{(\lambda,p)\in\Lambda_N^{\mathrm{phys}}}
\inf\!\left\{x\ge0:
\Pr\!\left[d_{H_N}(G_N+\lambda q_N(p),\mathcal C_N)^2\le x\right]
\ge1-\alpha\right\},
\]
and set
\[
C_N^{\mathrm{att}}=\{b\in\mathcal B:T_N(b)\le c_{N,\alpha}^{\mathrm{att}}\}.
\]
The handle uses the bounded-amplitude, tangent-plane, and endpoint-half-plane compactification to justify attainment and uniform transfer to the actual experiment.  It defines an exact mathematical critical value; a terminating certified numerical approximation is not part of the construction.
%%% FIELD calibration-handle-replacement
For each \(N\) and observed design, form the physically attainable set \(\Lambda_N^{\mathrm{phys}}\), define
\[
c_{N,\alpha}^{\mathrm{att}}
=\max_{(\lambda,p)\in\Lambda_N^{\mathrm{phys}}}
\inf\!\left\{x\ge0:
\Pr\!\left[d_{H_N}(G_N+\lambda q_N(p),\mathcal C_N)^2\le x\right]
\ge1-\alpha\right\},
\]
and set
\[
C_N^{\mathrm{att}}=\{b\in\mathcal B:T_N(b)\le c_{N,\alpha}^{\mathrm{att}}\}.
\]
The handle uses the bounded-amplitude, tangent-plane, and endpoint-half-plane compactification to justify attainment and uniform transfer to the actual experiment.  It defines an exact mathematical critical value; a terminating certified numerical approximation is not part of the construction.
%%% FIELD certified-computation-partial
The exact maximum \(c_{N,\alpha}^{\mathrm{att}}\) is attained, its reference quantile map is uniformly continuous on the proved compactification, and
\[
0<c_0\le c_{N,\alpha}^{\mathrm{att}}\le\chi^2_{2,1-\alpha}
\]
uniformly.  For each fixed input pair \((\lambda,p)\), the defining three-dimensional Gaussian probability is explicit.  These facts establish mathematical existence and a conservative analytic upper bound, but not a terminating certified approximation of the continuum maximum.
